% Options for packages loaded elsewhere
\PassOptionsToPackage{unicode}{hyperref}
\PassOptionsToPackage{hyphens}{url}
\documentclass[
]{article}
\usepackage{xcolor}
\usepackage[margin=1in]{geometry}
\usepackage{amsmath,amssymb}
\setcounter{secnumdepth}{-\maxdimen} % remove section numbering
\usepackage{iftex}
\ifPDFTeX
  \usepackage[T1]{fontenc}
  \usepackage[utf8]{inputenc}
  \usepackage{textcomp} % provide euro and other symbols
\else % if luatex or xetex
  \usepackage{unicode-math} % this also loads fontspec
  \defaultfontfeatures{Scale=MatchLowercase}
  \defaultfontfeatures[\rmfamily]{Ligatures=TeX,Scale=1}
\fi
\usepackage{lmodern}
\ifPDFTeX\else
  % xetex/luatex font selection
\fi
% Use upquote if available, for straight quotes in verbatim environments
\IfFileExists{upquote.sty}{\usepackage{upquote}}{}
\IfFileExists{microtype.sty}{% use microtype if available
  \usepackage[]{microtype}
  \UseMicrotypeSet[protrusion]{basicmath} % disable protrusion for tt fonts
}{}
\makeatletter
\@ifundefined{KOMAClassName}{% if non-KOMA class
  \IfFileExists{parskip.sty}{%
    \usepackage{parskip}
  }{% else
    \setlength{\parindent}{0pt}
    \setlength{\parskip}{6pt plus 2pt minus 1pt}}
}{% if KOMA class
  \KOMAoptions{parskip=half}}
\makeatother
\usepackage{graphicx}
\makeatletter
\newsavebox\pandoc@box
\newcommand*\pandocbounded[1]{% scales image to fit in text height/width
  \sbox\pandoc@box{#1}%
  \Gscale@div\@tempa{\textheight}{\dimexpr\ht\pandoc@box+\dp\pandoc@box\relax}%
  \Gscale@div\@tempb{\linewidth}{\wd\pandoc@box}%
  \ifdim\@tempb\p@<\@tempa\p@\let\@tempa\@tempb\fi% select the smaller of both
  \ifdim\@tempa\p@<\p@\scalebox{\@tempa}{\usebox\pandoc@box}%
  \else\usebox{\pandoc@box}%
  \fi%
}
% Set default figure placement to htbp
\def\fps@figure{htbp}
\makeatother
\setlength{\emergencystretch}{3em} % prevent overfull lines
\providecommand{\tightlist}{%
  \setlength{\itemsep}{0pt}\setlength{\parskip}{0pt}}
\usepackage{bookmark}
\IfFileExists{xurl.sty}{\usepackage{xurl}}{} % add URL line breaks if available
\urlstyle{same}
\hypersetup{
  pdftitle={NHIS 2021 Analysis},
  pdfauthor={Jordan Agosta, Lizbeth Ocampo, Lili Mora, Gabriela Quezada},
  hidelinks,
  pdfcreator={LaTeX via pandoc}}

\title{NHIS 2021 Analysis}
\author{Jordan Agosta, Lizbeth Ocampo, Lili Mora, Gabriela Quezada}
\date{2026-04-19}

\begin{document}
\maketitle

\section{Introduction}\label{introduction}

Physical health indicators such as body weight and height are widely
used in public health research to monitor population health trends.The
2021 National Health Interview Survey (NHIS) provides a nationally
representative sample of the U.S. population. This report presents a
descriptive analysis of self-reported health variables from the 2021
NHIS.

\section{Methods}\label{methods}

Data was drawn from the 2021 NHIS. The analytic sample consisted of
29,482 adult respondents after applying listwise deletion to remove
cases with missing or non-response codes per codebook specifications.
Variables of interest included age (AGEP\_A), height in inches
(HEIGHTTC\_A), weight in pounds (WEIGHTLBTC\_A), sex (SEX\_A), education
level (EDUCP\_A), and general health status (PHSTAT\_A). All analyses
were conducted in R using the tidyverse and psych packages.

\section{Results}\label{results}

Figure 1 displays the relationship between height and weight stratified
by sex and faceted by general health status. A positive association
between height and weight was observed across all health categories.
Males were consistently taller and heavier than females. Respondents
reporting Fair or Poor health showed greater weight variability compared
to those reporting Excellent or Very Good health.

\begin{figure}
\centering
\pandocbounded{\includegraphics[keepaspectratio]{Final_Report_files/Final_Report_files/figure-latex/plot1-1.pdf}}
\caption{Figure 1. Height vs.~weight by sex and general health status.}
\end{figure}

Figure 2 presents the Pearson correlation matrix for age, weight, and
height. The correlation between height and weight was r = 0.50,
indicating a moderate positive relationship. The correlations between
age and weight (r = -0.02) and age and height (r = -0.09) were both
weak.

\begin{figure}
\centering
\pandocbounded{\includegraphics[keepaspectratio]{Final_Report_files/Final_Report_files/figure-latex/plot2-1.pdf}}
\caption{Figure 2. Correlation matrix: age, weight, and height.}
\end{figure}

\section{Discussion}\label{discussion}

The positive height--weight relationship was consistent across all
health status groups, with greater dispersion observed in the Fair and
Poor categories. This may reflect the role of obesity-related chronic
conditions in poorer health outcomes. The weak correlations between age
and both height and weight are expected, as height is largely stable in
adulthood and weight changes variably throughout life.

\section{Conclusion}\label{conclusion}

This analysis of the 2021 NHIS provides a descriptive characterization
of the relationships among height, weight, age, sex, and general health
status in a nationally representative sample of U.S. adults. These
findings establish a foundation for future multivariate analyses
examining determinants of physical health after adjusting for
confounders such as race/ethnicity, education, and income.

\end{document}
